% Transcription — fonds Alexandre Grothendieck, shelfmark 129, batch 3 (pages 41-60).
% Established from the facsimile archives/batches/129/batch-03.pdf, prepared
% from a copy of 129.pdf fetched through the project's relay
% (https://grothendieck-relay.onrender.com/source/129.pdf; PDF page k =
% archivists' page 40+k, no cover sheet), per the skill transcribe-grothendieck.
% The folder is 114 pages in six batches, transcribed in parallel by six
% passes; this is the third (pages 41-60).
%
% Pass: Opus 5.5 (claude-opus-5-5), 2026-09-26 — first pass, unchecked against the pages by a human.
%
% Dating: \dating{[1970]} is the folder's own date in src/content/catalogue.ts,
% copied verbatim; since the folder has a date of its own, the group rule
% (« Descentes », 126-133) does not apply. Nothing on these leaves dates them.
%
% Pages skipped: none.
%
% Pages identified, not transcribed (other people's typescript): all twenty.
% By the folder's rule the typed text is never re-set; each page gets one
% short French \note identifying it (subject, section numbers), followed by
% every handwritten intervention on it, given in full. No author's name
% appears on these pages.
%
% Handwriting. Folios « II.13 » ... « II.26 » (pages 41-54, ink or pencil)
% and « III-1 » ... « III-5 » (pages 56-60, pencil), top right; page 55
% (chapter III title page) has none. Proof corrections of the same kind as
% batch 2's, on pages 41-43, 45-47, 50, 51, 53, 54 and 59: a cross (×, or
% « XX ») in the left margin, often with the corrected symbol repeated
% there, and the correction in the line (⊗ inked, ↦ inked, « objectif »
% struck and « bijectif » written, a surplus parenthesis and brace struck,
% λ and x inked, φ reinforced, a typed J replaced by j), plus two insertions
% in cursive: « conduisant aux identifications » (page 42, completing the
% sentence broken off at the foot of page 41) and « les » with an added
% plural s (page 43). Whether this hand is Grothendieck's cannot be told
% from the leaves: nothing in these corrections is personal to him, and the
% notes say so rather than attribute them; hence no \marginal{} anywhere.
% The script « ℓ.f. », the hollow letters and the script D of the typescript
% are uniform on every page, without margin crosses, and are taken as part
% of the typing, not as interventions. One faint mark before the folio of
% page 59 is \ill{}.
%
% His pagination: none of his. The folios are handwritten, hand undetermined.
%
% What the batch is about. (1) Pages 41-54: the rest of « Chapitre II.
% Groupes finis localement libres et théorie classique de Dieudonné »,
% continuing §4 from page 40 — the Dieudonné ring relations, the category
% of Dieudonné modules of finite length, Theorem 4.2 (Dieudonné: the
% anti-equivalence D* between finite commutative p-groups over a perfect
% field k and Dieudonné modules of finite length) and Corollary 4.2.1; 4.3
% the unipotent case (D*(G) = Hom(G, Witt ind-scheme), 4.3.2 compatibility
% with Frobenius and Verschiebung, Corollary 4.3.3); 4.4 the multiplicative
% case via Cartier and Pontryagin duality; 4.5 the general case and the
% fitting-type lemma M = M' × M''; §5 corollaries (5.1 base change to a
% perfect extension, after Oda and Oort; 5.2 rank = p^{length}; 5.3
% remarks on Barsotti-Tate groups and duality); §6 an annex constructing
% the quasi-inverse functor E(M) (6.1-6.4: representability by an affine
% scheme, E(M) finite unipotent, E adjoint to D* via Prim(G)). (2) Page 55:
% title page « Chapitre III. Groupes de Barsotti-Tate ». (3) Pages 56-60:
% chapter III, §1 notations (topology between fppf and fpqc, G(n), G(∞)),
% §2 flatness over Λ_n = Z/p^n Z and the criterion Prop. 2.2 (a)-(g) with
% its proof, 2.3, Prop. 2.4, and the start of §3 « Groupes de
% Barsotti-Tate tronqués d'échelon n », 3.1, breaking off where the batch
% ends.
%
% Hardest pages: 43 (the added s and « les ») and 53 (which strokes are
% deletions); page 59's faint mark near the folio. \ill{}: 1. \uncertain{}: 0.

\documentclass[11pt,a4paper]{article}
% Le préambule commun à toutes les transcriptions du fonds Grothendieck.
%
% Il tient en une idée : **l'incertitude se compose, elle ne se résout pas.**
% Une transcription qui lisse un mot douteux a détruit la seule chose pour
% laquelle on la faisait — un lecteur doit pouvoir savoir, à chaque endroit,
% ce qui était sur le feuillet et ce qui a été ajouté. Les six macros qui
% suivent sont donc l'appareil critique, et non de la décoration.
%
% Les mêmes commandes sont reconnues par `scripts/render.mjs`, qui produit la
% vue de lecture du site. Une macro ajoutée ici doit l'être aussi là-bas,
% faute de quoi le rendu échouera bruyamment — ce qui est voulu : un rendu
% silencieusement faux est pire qu'un rendu absent.

\ProvidesPackage{grothendieck}[2026/08/08 Transcriptions du fonds Grothendieck]

\usepackage{amsmath,amssymb,amsthm}
\usepackage{mathrsfs}
\usepackage{stmaryrd}
% Les diagrammes commutatifs sont partout dans ce fonds : c'est la langue
% même de la géométrie algébrique de Grothendieck, et une transcription qui
% les rendrait en prose ne transcrirait rien.
\usepackage{tikz-cd}
% Français par défaut — c'est la langue des feuillets — mais l'anglais est
% chargé aussi : la traduction et le résumé sont des documents anglais, et la
% typographie française (espaces avant les deux-points) y serait une faute.
% Ces documents basculent par \selectlanguage{english} en tête de corps.
\usepackage[english,french]{babel}
\usepackage{fontspec}
\usepackage{geometry}
\geometry{a4paper,margin=25mm,marginparwidth=28mm}
\usepackage{marginnote}
\usepackage{xcolor}
% ulem, not soul. `soul`'s \st and \ul are built on a character-by-character
% scan that XeLaTeX's font machinery breaks: the document fails with an
% undefined control sequence rather than degrading. ulem does the same two jobs
% and is Unicode-engine safe. `normalem` stops it hijacking \emph, which the
% transcriptions use for Grothendieck's own emphasis.
\usepackage[normalem]{ulem}
\usepackage{hyperref}
\hypersetup{colorlinks=true,linkcolor=black,urlcolor=black,pdfborder={0 0 0}}

% --- Métadonnées du lot -----------------------------------------------------
% Reprises telles quelles par le script de rendu, qui les affiche en tête de
% la vue de lecture. Elles fixent ce que le fichier prétend couvrir : sans
% elles, une transcription flotte sans point d'attache dans un fonds de seize
% mille feuillets.

\newcommand{\folder}[1]{\gdef\grothFolder{#1}}
\newcommand{\batch}[1]{\gdef\grothBatch{#1}}
\newcommand{\leaves}[2]{\gdef\grothFirst{#1}\gdef\grothLast{#2}}
\newcommand{\foldertitle}[1]{\gdef\grothFoldertitle{#1}}
\gdef\grothFolder{?}\gdef\grothBatch{?}\gdef\grothFirst{?}\gdef\grothLast{?}\gdef\grothFoldertitle{}

% --- Le résumé --------------------------------------------------------------
%
% Il ouvre la lecture modernisée, sous le titre « Résumé », et il est composé
% en petit. Ce n'est pas un ornement : c'est le seul passage du document qui
% ne soit pas des mathématiques, et un lecteur doit pouvoir le reconnaître
% avant de l'avoir lu — puis le sauter s'il connaît déjà le sujet, ou s'y
% arrêter s'il ne le connaît pas. Le filet qui le suit marque la reprise du
% corps.

\newenvironment{resume}
  {\small\setlength{\parskip}{0.45em}}
  {\par\medskip\hrule\medskip}

% --- Mots-clés ---------------------------------------------------------------
%
% En anglais, en clôture du résumé de la lecture modernisée : le vocabulaire
% moderne sous lequel chercher ce que les feuillets construisent. Le manifeste
% du site (`npm run manifest`) les extrait pour étiqueter la cote entière —
% cette ligne est la source unique des tags, il n'y a pas de fichier à côté.
\newcommand{\keywords}[1]{\par\smallskip\noindent
  {\footnotesize\textsc{Keywords} --- \textit{#1}}\par}

% --- L'appareil critique ----------------------------------------------------

% Le numéro de feuillet, en marge. C'est l'ancre qui permet de régler un
% désaccord sur une formule en regardant *un* feuillet plutôt qu'une chemise
% de six cents. Le site s'en sert aussi pour tourner les pages du fac-similé
% quand on fait défiler la transcription.
\newcommand{\leaf}[1]{%
  \marginnote{\footnotesize\color{gray!70}\textsf{#1}}%
  \label{leaf:#1}\ignorespaces}

% Illisible : on ne devine pas. Un mot inventé qui se lit comme les autres est
% le pire résultat possible.
\newcommand{\ill}{\textcolor{red!60!black}{[\,\ldots\,]}}

% Lecture douteuse : le mot est donné, et signalé comme incertain.
\newcommand{\uncertain}[1]{\uline{#1}}

% Ajout de l'éditeur — une lettre manquante, un mot sous-entendu.
\newcommand{\add}[1]{\textcolor{blue!50!black}{[#1]}}

% Biffé par Grothendieck lui-même. Ce qu'il a rayé fait partie du document :
% une rature dit souvent où la pensée a changé de direction.
\newcommand{\struck}[1]{\sout{#1}}

% Note du transcripteur, sur l'état matériel du feuillet.
\newcommand{\note}[1]{\par\smallskip{\small\color{gray!60!black}\textit{[#1]}}\par\smallskip}

% Note marginale de Grothendieck, distincte du corps du texte.
\newcommand{\marginal}[1]{\par\smallskip\noindent\hspace{1em}%
  {\small\color{gray!60!black}#1}\par\smallskip}

% --- Titre ------------------------------------------------------------------

\AtBeginDocument{%
  \begin{center}
    {\large\bfseries Fonds Alexandre Grothendieck --- cote n\textsuperscript{o}~\grothFolder}\\[2pt]
    {\itshape\grothFoldertitle}\\[4pt]
    {\small feuillets \grothFirst--\grothLast\ (lot \grothBatch)}\\[2pt]
    {\footnotesize\color{gray!60!black}Transcription établie d'après le fac-similé numérisé
     par l'Université de Montpellier.\\
     Les lectures incertaines sont soulignées, les illisibles notées [\,\ldots\,],
     les ajouts entre crochets.}
  \end{center}
  \vspace{1em}}

\endinput


\folder{129}
\batch{3}
\pages{41}{60}
\dating{[1970]}
\watermark{Édition de démonstration}
\foldertitle{Notes séminaire Montréal (Delale, Labute, Hakim, Messing) : copies de tapuscrits annotés (s.d.), notes manuscrites (s.d.)}

\begin{document}

\page{41}
\note{tapuscrit, chapitre~II, suite du n°~4.1 commencé page 40 : relations
de l'anneau de Dieudonné ($F\lambda = \lambda^{\sigma}F$,
$\lambda V = V\lambda^{\sigma}$, $FV = VF = p$), catégorie des modules de
Dieudonné de longueur finie, objets $(M, F_M, V_M)$ avec
$M^{\sigma} = M \otimes_W (W,\sigma)$, et énoncé du Théorème~4.2
(Dieudonné) ; la page s'arrête en cours de phrase. Texte d'autrui, non
transcrit. Interventions manuscrites, d'une main qu'on ne peut pas
identifier d'après le feuillet (rien ne permet de dire si c'est celle de
Grothendieck) : folio «~II.13~» ; en marge une croix et le signe
$\otimes$, devant $M^{\sigma} = M \otimes_W (W,\sigma)$, dont le
$\otimes$ dactylographié est repassé à l'encre}

\page{42}
\note{même tapuscrit, fin de l'énoncé du Théorème~4.2
($D^*(F_{G/k}) = F_{D^*(G)}$, $D^*(V_G) = V_{D^*(G)}$), Corollaire~4.2.1
(caractérisation des groupes infinitésimaux, unipotents,
bi-infinitésimaux, étales, de type multiplicatif par $F$ et $V$), renvoi à
Gabriel et Demazure, \emph{Groupes algébriques}, réduction aux cas
unipotent et multiplicatif (d'après 3.2), et n°~4.3 « Cas des groupes
unipotents », 4.3.1 (définition de $D^*(G)$, renvois à I,~3 et I,~4.2).
Interventions manuscrites, main non identifiée : folio «~II.14~» ; en
tête de page, au-dessus des formules dactylographiées, une ligne écrite à
l'encre qui achève la phrase restée ouverte au bas de la page 41 :
« conduisant aux identifications » ; en marge, une croix}

\page{43}
\note{même tapuscrit, action de $F$ et $V$ sur $D^*(G)$, structure de
$\mathcal{D}$-module, et n°~4.3.2 (isomorphisme
$D^*(G^{(p)}) \simeq (D^*(G))^{\sigma}$ compatible au Frobenius et à la
Verschiebung ; renvoi à I,~5.2). Interventions manuscrites, main non
identifiée : folio «~II.15~» ; en marge «~XX~», devant « des morphismes
correspondant sur $W_n$ » : un s ajouté à la main à « correspondant », et
« les » inséré par un signe d'insertion avant $W_n$, ce qui donne « des
morphismes correspondants sur les $W_n$ » ; plus bas, en marge une croix
et le signe $\mapsto$, devant la ligne « $u \mapsto u \circ F_{G/k}$ », dont
le $\mapsto$ est tracé à l'encre}

\page{44}
\note{même tapuscrit, fin du n°~4.3.2 (renvoi à I,~5.3), cas unipotent
(renvoi à 2.5.2) et Corollaire~4.3.3 (trois équivalences de catégories :
$p$-groupes finis unipotents, étales, bi-infinitésimaux, et les
sous-catégories correspondantes de modules de Dieudonné ; renvoi à 2.5).
Seule intervention manuscrite : le folio «~II.16~»}

\page{45}
\note{même tapuscrit, n°~4.4 « Cas des groupes de type multiplicatif » :
réduction, par la dualité de Cartier et le Corollaire~4.3.3, aux
$p$-groupes étales, et dualité de Pontryagin
$G \mapsto \underline{\mathrm{Hom}}_{\mathrm{gr}}(G,
\mathbb{Q}_p/\mathbb{Z}_p)$. Interventions manuscrites, main non
identifiée : folio «~II.17~» ; en marge, devant cette dernière formule, une
croix et $\mathbb{Q}_p/\mathbb{Z}_p$, sans correction visible dans la
ligne}

\page{46}
\note{même tapuscrit, diagramme résumant le cas multiplicatif (dualités de
Cartier et de Pontryagin, renvoi à 4.3.3) et n°~4.5 « Cas général »
($G = G^{\mathrm{uni}} \times G^{\mathrm{mult}}$ d'après 3.2,
$D^*(G) = D^*(G^{\mathrm{uni}}) \times D^*(G^{\mathrm{mult}})$). Interventions
manuscrites, main non identifiée : folio «~II.18~» ; en marge une croix, et
dans la première ligne du diagramme (« formée des objets où $V$ est
objectif ») le mot « objectif » biffé d'un trait oblique, « bijectif »
écrit au-dessous}

\page{47}
\note{même tapuscrit, équivalence de la catégorie des modules de Dieudonné
avec un produit de deux sous-catégories ($V$ nilpotent, $V$ bijectif),
Lemme (décomposition $M = M' \times M''$) et sa preuve, puis §5
« Corollaires et compléments », Corollaire~5.1 (changement de base à une
extension parfaite $K$ de $k$). Interventions manuscrites, main non
identifiée : folio «~II.19~» ; en marge une croix, et dans la formule
détachée du haut la flèche marquée $\approx$ est tracée à l'encre}

\page{48}
\note{même tapuscrit, preuve du Corollaire~5.1 (d'après Oda et Oort),
Corollaire~5.2 (rang de $G$ égal à $p$ puissance la longueur de $D^*(G)$)
et sa preuve (cas de $\alpha_p$ et de $\mathbb{Z}/p\mathbb{Z}$, renvoi à
2.4). Seule intervention manuscrite : le folio «~II.20~», au crayon}

\page{49}
\note{même tapuscrit, n°~5.3 « Remarques » : 5.3.1 (groupes formels et
groupes de Barsotti-Tate, renvoi à III,~5.6), 5.3.2 (dualité
$D^*(G^*) \simeq \mathrm{Hom}_W(D^*(G), K/W)$, renvoi à Gabriel et
Demazure) ; début du §6 « Annexe~= Construction du foncteur quasi inverse
de $D^*$ ». Seule intervention manuscrite : le folio «~II.21~»}

\page{50}
\note{même tapuscrit, suite du §6 : n°~6.1 (le foncteur en groupes
$\mathbb{E}(M)$) et n°~6.2 « $\mathbb{E}(M)$ est représentable par un
schéma affine » (polynômes $P_N$, $S_N$, relations i)--iii) définissant la
$k$-algèbre $A$). Interventions manuscrites, main non identifiée : folio
«~II.22~» ; en marge une croix, devant la relation iii), où les quatre
$\lambda$ (dans $T_{\lambda x}$ et dans les arguments de $P_N$) sont écrits
à la main, et où les deux derniers $x$ (dans $T_{V^N(x)}$ et $T_x$) sont
repassés à l'encre}

\page{51}
\note{même tapuscrit, suite du n°~6.2 (factorisation par
$\mathbb{W}_{N+1}(S)$, isomorphisme fonctoriel $\phi_S$ et son inverse) et
début du n°~6.3 « Le schéma $\mathbb{E}(M)$ est un $p$-groupe fini
unipotent ». Interventions manuscrites, main non identifiée : folio
«~II.23~» ; en marge une croix, devant la ligne « Il est clair que
$\phi_S(u)$ est bien défini », dont le $\phi$ est tracé ou repassé à
l'encre}

\page{52}
\note{même tapuscrit, suite du n°~6.3 : $\mathbb{E}(M)$ est fini (système
de générateurs), puis unipotent (renvois à 2.5.2 et I.5.3 ; diagramme
commutatif reliant $\phi_S$ et $\phi_{S,\sigma}$, morphisme identifié à la
Verschiebung). Seule intervention manuscrite : le folio «~II.24~»}

\page{53}
\note{même tapuscrit, n°~6.4 « Le foncteur $\mathbb{E}$ est adjoint à
gauche de $D^*$ » : définition de $\mathrm{Prim}(G)$, isomorphismes
$D^*(G) \simeq \mathrm{Prim}(G)$ et
$\mathrm{Hom}_k(G, \mathbb{E}(M)) = \mathrm{Hom}_{\mathcal{D}}(M,
\mathbb{W}(G))$, diagramme commutatif pour $\Delta$. Interventions
manuscrites, main non identifiée : folio «~II.25~» ; en marge «~XX~»,
devant la formule qui définit $\mathrm{Prim}(G)$, où deux traits obliques
biffent une parenthèse ouvrante en trop après « Prim » et l'accolade
fermante en trop à la fin de la formule}

\page{54}
\note{même tapuscrit, fin du n°~6.4 : calcul de $\phi_S(u+u')$ et
condition pour que $\bar u$ soit un morphisme de groupes, d'où le noyau
contenu dans $\mathrm{Prim}(G)$ ; la page s'arrête à mi-hauteur, fin du
chapitre~II. Interventions manuscrites, main non identifiée : folio
«~II.26~» ; en marge une croix et le signe $\otimes$, devant
« $\mathbb{W}_{N+1}(\Gamma(\mathcal{O}_G) \otimes
\Gamma(\mathcal{O}_G))$ », où le signe dactylographié entre les deux
facteurs est entouré à l'encre et souligné d'un trait}

\page{55}
\note{page de titre dactylographiée : « Chapitre III. Groupes de
Barsotti-Tate ». Aucun nom d'auteur ; aucune intervention manuscrite}

\page{56}
\note{tapuscrit, chapitre~III « Groupes de Barsotti-Tate », §1
« Notations et définitions préliminaires » : topologie sur Sch/S entre
f.p.p.f. et f.p.q.c. (renvoi à SGA~3, IV, 6.3), groupes comme faisceaux en
groupes commutatifs, $\Lambda_n = \mathbb{Z}/p^n\mathbb{Z}$, $G(n)$,
$G(\infty)$ (renvoi à 5.1). Texte d'autrui, non transcrit. Seule
intervention manuscrite, main non identifiée : le folio «~III-1~», au
crayon}

\page{57}
\note{même tapuscrit, §2 « Platitude et critère de représentabilité » :
n°~2.1 (platitude sur $\Lambda_n$), n°~2.2 Proposition (équivalences (a)
à (g) pour un groupe $G(n)$ annulé par $p^n$) et début de la
démonstration. Seule intervention manuscrite : le folio «~III-2~», au
crayon}

\page{58}
\note{même tapuscrit, suite de la démonstration de 2.2 : (b)
$\Leftrightarrow$ (c) par des triangles commutatifs, (c) $\Rightarrow$
(a), (d) $\Rightarrow$ (e) par un diagramme à lignes exactes, début de
(e) $\Leftrightarrow$ (f). Seule intervention manuscrite : le folio
«~III-3~», au crayon}

\page{59}
\note{même tapuscrit, fin de la démonstration de 2.2 ((e) $\Leftrightarrow$
(f), (e) $\Rightarrow$ (g) $\Rightarrow$ (c)) et début du n°~2.3
(représentabilité de $G(n)$ par un schéma fini localement libre).
Interventions manuscrites, main non identifiée : folio «~III-4~», au
crayon, précédé d'une marque pâle \ill{} ; en marge une croix et $j$,
devant la formule $\beta^i = j \circ \lambda^{n-2} \circ \dots \circ
\lambda^i$, où le J majuscule dactylographié est biffé d'un trait oblique
et un $j$ cursif écrit au-dessus}

\page{60}
\note{même tapuscrit, fin du n°~2.3 (critère de platitude fibre à fibre,
renvoi à EGA~IV, 11.3.1 ; renvoi à II,~1.6), n°~2.4 Proposition, et début
du §3 « Groupes de Barsotti-Tate tronqués d'échelon $n$ », n°~3.1
(condition supplémentaire pour $n = 1$, renvoi à 3.3 ; morphismes
$F_{G_0/S_0}$ et $V_{G_0}$) ; la page s'arrête où finit le lot. Seule
intervention manuscrite : le folio «~III-5~», au crayon}

\end{document}
